\documentclass[11pt,a4paper]{article}

% ---------------------------------------------------------
% PREAMBLE (stable, warning-free, Overleaf-safe)
% ---------------------------------------------------------
\usepackage[utf8]{inputenc}
\usepackage[T1]{fontenc}
\usepackage{lmodern}
\usepackage{amsmath, amssymb, amsfonts}
\usepackage{bm}
\usepackage{geometry}
\usepackage{graphicx}
\usepackage{hyperref}
\usepackage{cite}
\usepackage{authblk}
\usepackage{physics}
\usepackage{mathtools}
\usepackage{textgreek}
\geometry{margin=2.4cm}

\title{\textbf{Companion LXXIII: Maximum Mean Discrepancy of Persistence Diagrams in the Relaxation-Driven Cyclic Cosmology}}
\author{Michael Lehmann \\ \texttt{mi.lehmann@gmx.de}}
\date{April 2026}

\begin{document}
\maketitle

\begin{abstract}
The maximum mean discrepancy (MMD) provides a kernel-based statistical test for 
distinguishing probability distributions in a reproducing kernel Hilbert space 
(RKHS). Applied to persistence diagrams, MMD quantifies whether two ensembles of 
topological features arise from the same underlying cosmological model. In weak-
lensing analyses, this enables a direct comparison between the Relaxation-Driven 
Cyclic Cosmology (RDCC) and the standard ΛCDM model. Building on the Hilbert-space 
embeddings and kernel PCA developed in previous Companions, this work derives the 
MMD for persistence diagrams in RDCC, identifies the scale-dependent signatures of 
potential relaxation, and establishes a statistically rigorous framework for 
model discrimination using topological data analysis. The resulting test provides 
a robust, nonparametric, and observationally accessible probe of the RDCC 
gravitational field.
\end{abstract}
% ---------------------------------------------------------
\section{Introduction}

Persistence diagrams encode the birth and death of topological features in weak-
lensing convergence fields. Hilbert-space embeddings and kernel PCA provide a 
linearized and spectrally decomposed representation of these diagrams. Maximum 
mean discrepancy (MMD) extends this framework by enabling statistical tests 
between two distributions of embedded diagrams.

In cosmology, MMD offers a powerful tool for distinguishing between models based 
on their topological signatures. In RDCC, the relaxation mechanism modifies the 
scale-dependent evolution of the gravitational potential, altering the distribution 
of persistence diagrams in a characteristic way. MMD provides a direct, 
nonparametric test for detecting these modifications in weak-lensing surveys.
% ---------------------------------------------------------
\section{Maximum Mean Discrepancy for Persistence Diagrams}

Given two ensembles of persistence diagrams


\[
    \{D_{i}\}_{i=1}^{n}, \qquad
    \{E_{j}\}_{j=1}^{m},
\]


their embeddings into an RKHS $\mathcal{H}$ are


\[
    \Phi(D_{i}), \qquad \Phi(E_{j}).
\]



The squared MMD is defined as
\begin{align}
    \mathrm{MMD}^{2}
    &= \left\|
        \frac{1}{n} \sum_{i=1}^{n} \Phi(D_{i})
        - \frac{1}{m} \sum_{j=1}^{m} \Phi(E_{j})
       \right\|_{\mathcal{H}}^{2} \\
    &= \frac{1}{n^{2}} \sum_{i,i'} k(D_{i}, D_{i'})
     + \frac{1}{m^{2}} \sum_{j,j'} k(E_{j}, E_{j'})
     - \frac{2}{nm} \sum_{i,j} k(D_{i}, E_{j}).
\end{align}

This quantity is zero if and only if the two distributions are identical.
% ---------------------------------------------------------
\section{RDCC Modification of Persistence Diagram Distributions}

In RDCC, the convergence evolves as
\begin{equation}
    \kappa_{\mathrm{RDCC}}(\ell)
    = \kappa(\ell)
      \left[
        1 - \chi_{\kappa}
        \left(\frac{\ell}{\ell_{\mathrm{rel}}}\right)^{\alpha}
      \right].
\end{equation}

This modifies the birth and death thresholds of topological features, shifting the 
distribution of persistence diagrams:


\[
    D \mapsto D_{\mathrm{RDCC}}.
\]



The kernel becomes


\[
    k_{\mathrm{RDCC}}(D_{i}, D_{j})
    = k(D_{i}, D_{j})
      \left[
        1 - \chi_{k}
        \left(\frac{\ell}{\ell_{\mathrm{rel}}}\right)^{\alpha}
      \right].
\]


% ---------------------------------------------------------
\section{MMD in RDCC}

The RDCC-modified MMD is

\begin{align}
    \mathrm{MMD}^{2}_{\mathrm{RDCC}}
    &= \mathrm{MMD}^{2}
       \left[
         1 - \chi_{\mathrm{MMD}}
         \left(\frac{\ell}{\ell_{\mathrm{rel}}}\right)^{\alpha}
       \right].
\end{align}

RDCC therefore induces:

\begin{itemize}
    \item increased separation from ΛCDM at large scales,
    \item suppressed separation at small scales,
    \item a characteristic turnover near $\ell_{\mathrm{rel}}$,
    \item enhanced sensitivity to long-lived persistence features,
    \item modified kernel expectations across redshift.
\end{itemize}
% ---------------------------------------------------------
\section{Hypothesis Testing: RDCC vs. ΛCDM}

Define the null hypothesis


\[
    H_{0}: \mathcal{P}_{\mathrm{RDCC}} = \mathcal{P}_{\Lambda\mathrm{CDM}},
\]


and the alternative


\[
    H_{1}: \mathcal{P}_{\mathrm{RDCC}} \neq \mathcal{P}_{\Lambda\mathrm{CDM}}.
\]



The test statistic is $\mathrm{MMD}^{2}_{\mathrm{RDCC}}$.

Permutation tests or asymptotic bounds yield $p$-values.

RDCC predicts:

\begin{itemize}
    \item significant rejection of $H_{0}$ at large scales,
    \item marginal separation at small scales,
    \item maximal discriminative power near $\ell_{\mathrm{rel}}$.
\end{itemize}
% ---------------------------------------------------------
\section{Conclusion}

Maximum mean discrepancy provides a statistically rigorous method for 
distinguishing the Relaxation-Driven Cyclic Cosmology from ΛCDM using persistence 
diagrams. The relaxation mechanism modifies the distribution of topological 
features in a scale-dependent manner, producing distinctive signatures in the MMD 
test statistic. These effects offer a direct, nonparametric, and observationally 
accessible probe of the RDCC gravitational field. The present Companion establishes 
the foundation for future studies of statistical topology, model discrimination, 
and cosmological inference in RDCC.

\bibliographystyle{apsrev4-2}
\nocite{*}
\bibliography{references}


\end{document}
